% Chapter 4, Topic _Linear Algebra_ Jim Hefferon
%  http://joshua.smcvt.edu/linearalgebra
%  2001-Jun-12
\topic{Cramer 法则}
\index{Cramer's rule|(}
% We have introduced determinant functions algebraically by looking
% for a formula to decide whether a matrix is nonsingular.
% After that introduction we saw a geometric interpretation, 
% that the determinant function
% gives the size of the box with sides formed by the columns of the matrix.
% Here we make a connection between the two views.

%<*CramersRuleExample0>
一个线性方程组等价于向量之间的一个线性关系。
\begin{equation*}
  \begin{linsys}{2}
     x_1  &+  &2x_2  &=  &6  \\
    3x_1  &+  &x_2   &=  &8 
  \end{linsys}
% \end{equation*}
% \begin{equation*}
  \qquad\Longleftrightarrow\qquad
  x_1\cdot\colvec{1 \\ 3}+x_2\cdot\colvec{2 \\ 1}=\colvec{6 \\ 8}
\end{equation*}
%</CramersRuleExample0>
在下图中，
小的平行四边形由向量
$\binom{1}{3}$ 与 $\binom{2}{1}$ 构成。
它嵌套在一个以 $x_1\binom{1}{3}$ 与~$x_2\binom{2}{1}$ 为边的平行四边形之内。
由向量方程可知，较大平行四边形的远端顶点是 $\binom{6}{8}$。
\begin{center}
  \includegraphics{det/mp/ch4.1}
\end{center}
%<*CramersRuleExample1>
这幅图把求线性方程组的解这一代数问题
改用几何语言来陈述：~我们
要把起始平行四边形的两条边各伸缩多少倍（即 $x_1$ 和~$x_2$ 取什么值），
才能使它填满另一个平行四边形？
%</CramersRuleExample1>

我们可以利用这幅图，以及对行列式的几何理解，
来得到求解线性方程组的一个新公式。
比较这些带阴影的箱体的大小。
\begin{center}
   \includegraphics{det/mp/ch4.2}
   \hfil
   \includegraphics{det/mp/ch4.3}
   \hfil
   \includegraphics{det/mp/ch4.4}
\end{center}
%<*CramersRuleExample2>
第二个箱体由向量 $x_1\binom{1}{3}$ 与 $\binom{2}{1}$ 确定，而
大小函数\Dash 即行列式\Dash 的性质之一是：
第二个箱体的大小是第一个的 \( x_1 \) 倍。
第三个箱体由第二个箱体经剪切而得：把
$x_2\binom{2}{1}$ 加到 $x_1\binom{1}{3}$ 上，得到
$x_1\binom{1}{3}+x_2\binom{2}{1}=\binom{6}{8}$，
同时保留 $\binom{2}{1}$。
行列式不受剪切的影响，所以
第三个箱体的大小等于第二个的大小。
%</CramersRuleExample2>

%<*CramersRuleExample3>
综合起来，我们得到下式。
\begin{equation*}
  x_1\cdot \begin{vmat}[r]
     1  &2  \\
     3  &1
  \end{vmat}
  =
  \begin{vmat}
     x_1\cdot 1  &2  \\
     x_1\cdot 3  &1
  \end{vmat}
  =
  \begin{vmat}
     x_1\cdot 1+x_2\cdot 2  &2  \\
     x_1\cdot 3+x_2\cdot 1  &1
  \end{vmat}
  =
  \begin{vmat}[r]
     6  &2  \\
     8  &1
  \end{vmat}
\end{equation*}
%</CramersRuleExample3>
%<*CramersRuleExample4>
解出其中一个变量的值。
\begin{equation*}
  x_1=
  \frac{\begin{vmat}[r]
     6  &2  \\
     8  &1
  \end{vmat} }{
  \begin{vmat}[r]
     1  &2  \\
     3  &1
  \end{vmat}  }
  =\frac{-10}{-5}=2
\end{equation*}
%</CramersRuleExample4>

这个例子的推广就是\definend{Cramer 法则}（克拉默法则）：%
\index{determinant!Cramer's rule}%
\index{linear equation!solution of!Cramer's rule}
若 \( \deter{A}\neq 0 \)，则方程组 \( A\vec{x}=\vec{b} \) 有唯一解
$
   x_i=\deter{B_i}/\deter{A}
$
其中矩阵 $B_i$ 是把 $A$ 的第~$i$ 列换成
向量 \( \vec{b} \) 而得到的。
其证明见 \nearbyexercise{ex:CramerRule}。

例如，要求出下方程组中的 \( x_2 \)
\begin{equation*}
  \begin{mat}[r]
    1  &0  &4  \\
    2  &1  &-1 \\
    1  &0  &1
  \end{mat}
  \colvec{x_1 \\ x_2 \\ x_3}
  =\colvec[r]{2 \\ 1 \\ -1}
\end{equation*}
我们做如下计算。
\begin{equation*}
  x_2=
  \frac{ \begin{vmat}[r]
           1  &2  &4  \\
           2  &1  &-1 \\
           1  &-1 &1
         \end{vmat}  }{
         \begin{vmat}[r]
           1  &0  &4  \\
           2  &1  &-1 \\
           1  &0  &1
         \end{vmat}  }
  =\frac{-18}{-3}
\end{equation*}

Cramer 法则让我们能够凭观察求解小而简单的方程组。
例如，对于两个方程两个未知数、或三个方程三个未知数、
且数字都是较小整数的方程组，我们可以直接求解。
这类情形出现得相当频繁，因此许多人觉得这个公式很顺手。

但无论是手工还是用计算机，用它来求解大型或复杂的方程组都不切实际。
基于高斯消元法的方法更快。

\begin{exercises}
  \item
    用 Cramer 法则求下列各方程组中每个变量的值。
    \begin{exparts*}
      \partsitem $\begin{linsys}{2}
                    x  &- &y  &=  &4  \\
                   -x  &+ &2y &=  &-7
                  \end{linsys}$
      \partsitem $\begin{linsys}{2}
                    -2x  &+  &y  &=  &-2 \\
                      x  &-  &2y &=  &-2
                  \end{linsys}$
    \end{exparts*}
    \begin{answer}
      \begin{exparts*}
        \partsitem 分别解出各个变量。
          \begin{equation*}
            x=
             \frac{ \begin{vmat}[r]
                       4  &-1  \\
                      -7  &2
                    \end{vmat}  }{
                    \begin{vmat}[r]
                       1  &-1  \\
                      -1  &2
                    \end{vmat}  }
            =\frac{1}{1}=1
            \qquad
            y=
             \frac{ \begin{vmat}[r]
                       1  &4  \\
                      -1  &-7
                    \end{vmat}  }{
                    \begin{vmat}[r]
                       1  &-1  \\
                      -1  &2
                    \end{vmat}  }
            =\frac{-3}{1}=-3
          \end{equation*}
        \partsitem $x=2$，$y=2$
      \end{exparts*}
    \end{answer}
  \item
    用 Cramer 法则求该方程组中的 \( z \)。
    \begin{equation*}
      \begin{linsys}{4}
        2x  &+  &y  &+  &z  &=  &1 \\
        3x  &   &   &+  &z  &=  &4 \\
         x  &-  &y  &-  &z  &=  &2
      \end{linsys}
    \end{equation*}
    \begin{answer}
      \( z=1 \)
    \end{answer}
  \item \label{ex:CramerRule}
    证明 Cramer 法则。
    \begin{answer}
      行列式在倍加下不变，
      包括列的倍加，于是
      \( \det(B_i)=\det(\vec{a}_1,\dots,
      x_1\vec{a}_1+\dots+x_i\vec{a}_i+\dots+x_n\vec{a}_n,\dots,\vec{a}_n) \)。
      用“取第一列的 $-x_1$ 倍加到第~$i$ 列”等操作，可以看出它
      等于
      $\det(\vec{a}_1,\dots,x_i\vec{a}_i,\dots,\vec{a}_n)$。
      进而，它又等于
      $x_i\cdot\det(\vec{a}_1,\dots,\vec{a}_i,\dots,\vec{a}_n)
         =x_i\cdot\det(A)$，
      这正是所要证的。
    \end{answer}
  \item
   下面是 Cramer 法则的另一种证明，其中不显含任何几何内容。
   记 $X_i$ 为把第~$i$ 列换成由未知数
   $x_1$、\ldots、~$x_n$ 组成的向量~$\vec{x}$
   后所得的单位矩阵。
   \begin{exparts}
     \partsitem 验证 $AX_i=B_i$。
     \partsitem 对两边取行列式。
   \end{exparts}
   \begin{answer}
     \begin{exparts*}
       \partsitem
         下面是 $\nbyn{2}$ 方程组取 $i=2$ 的情形。
         \begin{equation*}
           \begin{linsys}{2}
             a_{1,1}x_1 &+ &a_{1,2}x_2 &= &b_1 \\
             a_{2,1}x_1 &+ &a_{2,2}x_2 &= &b_2
           \end{linsys}
           \quad\Longleftrightarrow\quad
           \begin{mat}
             a_{1,1}  &a_{1,2}  \\
             a_{2,1}  &a_{2,2}
           \end{mat}
           \begin{mat}
             1  &x_1 \\
             0  &x_2
           \end{mat}
           =
           \begin{mat}
             a_{1,1}  &b_1 \\
             a_{2,1}  &b_2
           \end{mat}
         \end{equation*}
      \partsitem 行列式函数是乘性的：
        $\det(B_i)=\det(AX_i)=\det(A)\cdot\det(X_i)$。
        Laplace 展开表明 $\det(X_i)=x_i$，
        解出 $x_i$ 便得到 Cramer 法则。
     \end{exparts*}
   \end{answer}

  \item
    设一个线性方程组的方程个数与未知数个数相同，
    其所有系数与常数都是整数，且其系数
    矩阵的行列式为~\( 1 \)。
    证明解中的元素全为整数。
    （\textit{注记。}
     这常被用来为练习编造线性方程组。）
    \begin{answer}
      由于 $A$ 的行列式非零，Cramer 法则适用，并给出 $x_i=\deter{B_i}/1$。
      由于 $B_i$ 是整数矩阵，故其行列式是整数。
    \end{answer}
  \item
    用 Cramer 法则给出二个方程/二个未知数的线性方程组的解的公式。
    \begin{answer}
      方程组
      \begin{equation*}
        \begin{linsys}{2}
           ax  &+  by  &=  &e  \\
           cx  &+  dy  &=  &f
        \end{linsys}
      \end{equation*}
      的解为
      \begin{equation*}
        x=\frac{ed-fb}{ad-bc}
        \qquad
        y=\frac{af-ec}{ad-bc}
      \end{equation*}
      当然前提是分母不为零。
    \end{answer}
  \item
    Cramer 法则能区分无解的方程组与
    有无穷多解的方程组吗？
    \begin{answer}
      当然，奇异方程组的 \( \deter{A} \) 等于零；
      而无穷多解的情形可以用
      所有的 \( \deter{B_i} \) 也都为零这一事实来刻画。
    \end{answer}
  \item
    本专题的第一幅图（即不使用行列式的那幅）
    展示的是唯一解的情形。
    请对无穷多解的情形
    与无解的情形各画一幅类似的图。
    \begin{answer}
      我们可以把两种非奇异的情形与下面这个
      方程组一起考虑
      \begin{equation*}
        \begin{linsys}{2}
           x_1  &+  &2x_2  &=  &6  \\
           x_1  &+  &2x_2  &=  &c
        \end{linsys}
      \end{equation*}
      其中 $c=6$ 当然给出无穷多解，而 $c$ 取任何其他值
      都给出无解。
      相应的向量方程
      \begin{equation*}
        x_1\cdot\colvec[r]{1 \\ 1}+x_2\cdot\colvec[r]{2 \\ 2}=\colvec[r]{6 \\ c}
      \end{equation*}
      给出两条重叠向量的图景。
      它们都位于直线 $y=x$ 上。
      在 $c=6$ 的情形，右边的向量也位于
      直线 $y=x$ 上，而在其他任何情形下都不是。
    \end{answer}
\end{exercises}
\index{Cramer's rule|)}
\endinput
